\documentclass[a4paper,11pt]{book}
%\documentclass[a4paper,twoside,11pt,titlepage]{book}
\usepackage{listings}
\usepackage[utf8]{inputenc}
\usepackage[spanish]{babel}

% \usepackage[style=list, number=none]{glossary} %
%\usepackage{titlesec}
%\usepackage{pailatino}

\usepackage[
backend=biber,
style=alphabetic,
sorting=ynt
]{biblatex}
\addbibresource{bibliografia.bib}

\usepackage{breakcites}
\usepackage{float}

\decimalpoint
\usepackage{dcolumn}
\newcolumntype{.}{D{.}{\esperiod}{-1}}
\makeatletter
\addto\shorthandsspanish{\let\esperiod\es@period@code}
\makeatother


%\usepackage[chapter]{algorithm}
\RequirePackage{verbatim}
%\RequirePackage[Glenn]{fncychap}
\usepackage{fancyhdr}
\usepackage{graphicx}
\graphicspath{./doc/imagenes}
\usepackage{afterpage}

\usepackage{longtable}

\usepackage[pdfborder={000}]{hyperref} %referencia

% ********************************************************************
% Re-usable information
% ********************************************************************
\newcommand{\myTitle}{Aplicación de renderizado 3D\xspace}
\newcommand{\myDegree}{Grado en ingeniería informática\xspace}
\newcommand{\myName}{Pablo Cantudo Gómez\xspace}
\newcommand{\myProf}{Juan Carlos Torres Cantero y Luis López Escudero\xspace}
%\newcommand{\mySupervisor}{Put name here\xspace}
\newcommand{\myFaculty}{Escuela Técnica Superior de Ingenierías Informática y de
Telecomunicación\xspace}
\newcommand{\myFacultyShort}{E.T.S. de Ingenierías Informática y de
Telecomunicación\xspace}
\newcommand{\myDepartment}{Departamento de Lenguajes y Sistemas Informáticos\xspace}
\newcommand{\myUni}{\protect{Universidad de Granada}\xspace}
\newcommand{\myLocation}{Granada\xspace}
\newcommand{\myTime}{\today\xspace}
\newcommand{\myVersion}{Version 0.1\xspace}


\usepackage{dirtree}
\usepackage{url}
\usepackage{colortbl,longtable}
\usepackage[stable]{footmisc}
\usepackage{index}

\makeindex
\usepackage[toc,nonumberlist]{glossaries}
\makeglossaries

% Definición de comandos que me son tiles:
%\renewcommand{\indexname}{Índice alfabético}
%\renewcommand{\glossaryname}{Glosario}

\pagestyle{fancy}
\fancyhf{}
\fancyhead[LO]{\leftmark}
\fancyhead[RE]{\rightmark}
\fancyhead[RO,LE]{\textbf{\thepage}}
\renewcommand{\chaptermark}[1]{\markboth{\textbf{#1}}{}}
\renewcommand{\sectionmark}[1]{\markright{\textbf{\thesection. #1}}}

\setlength{\headheight}{1.5\headheight}

\newcommand{\HRule}{\rule{\linewidth}{0.5mm}}
%Definimos los tipos teorema, ejemplo y definición podremos usar estos tipos
%simplemente poniendo \begin{teorema} \end{teorema} ...
\newtheorem{teorema}{Teorema}[chapter]
\newtheorem{ejemplo}{Ejemplo}[chapter]
\newtheorem{definicion}{Definición}[chapter]

\definecolor{gray97}{gray}{.97}
\definecolor{gray75}{gray}{.75}
\definecolor{gray45}{gray}{.45}
\definecolor{gray30}{gray}{.94}

\lstset{ frame=Ltb,
     framerule=0.5pt,
     aboveskip=0.5cm,
     framextopmargin=3pt,
     framexbottommargin=3pt,
     framexleftmargin=0.1cm,
     framesep=0pt,
     rulesep=.4pt,
     backgroundcolor=\color{gray97},
     rulesepcolor=\color{black},
     %
     stringstyle=\ttfamily,
     showstringspaces = false,
     basicstyle=\scriptsize\ttfamily,
     commentstyle=\color{gray45},
     keywordstyle=\bfseries,
     %
     numbers=left,
     numbersep=6pt,
     numberstyle=\tiny,
     numberfirstline = false,
     breaklines=true,
   }
 
% minimizar fragmentado de listados
\lstnewenvironment{listing}[1][]
   {\lstset{#1}\pagebreak[0]}{\pagebreak[0]}

\lstdefinestyle{CodigoC}
   {
	basicstyle=\scriptsize,
	frame=single,
	language=C,
	numbers=left
   }
\lstdefinestyle{python}
   {
	basicstyle=\small,
	frame=single,
	backgroundcolor=\color{gray30},
	language=C++,
	numbers=left
   }

 
\lstdefinestyle{Consola}
   {basicstyle=\scriptsize\bf\ttfamily,
    backgroundcolor=\color{gray30},
    frame=single,
    numbers=none
   }


\newcommand{\bigrule}{\titlerule[0.5mm]}


%Para conseguir que en las páginas en blanco no ponga cabecerass
\makeatletter
\def\clearpage{%
  \ifvmode
    \ifnum \@dbltopnum =\m@ne
      \ifdim \pagetotal <\topskip
        \hbox{}
      \fi
    \fi
  \fi
  \newpage
  \thispagestyle{empty}
  \write\m@ne{}
  \vbox{}
  \penalty -\@Mi
}
\makeatother

\usepackage{pdfpages}
\begin{document}
\begin{titlepage}
 
 
\newlength{\centeroffset}
\setlength{\centeroffset}{-0.5\oddsidemargin}
\addtolength{\centeroffset}{0.5\evensidemargin}
\thispagestyle{empty}

\noindent\hspace*{\centeroffset}\begin{minipage}{\textwidth}

\centering
\includegraphics[width=0.9\textwidth]{imagenes/logo.png}\\[1.4cm]

\textsc{ \Large TRABAJO FIN DE GRADO\\[0.2cm]}
\textsc{ GRADO EN INGENIERÍA INFORMÁTICA}\\[1cm]
% Upper part of the page
% 
% Title
{\Huge\bfseries \\ Desarrollo de una aplicaci´on
de modelado 3D}
\noindent\rule[-1ex]{\textwidth}{3pt}\\[3.5ex]
{\large\bfseries Motor de renderizado en tiempo real basada en SDFs}
\end{minipage}

\vspace{2.5cm}
\noindent\hspace*{\centeroffset}\begin{minipage}{\textwidth}
\centering

\textbf{Autor}\\ {Pablo Cantudo Gómez}\\[2.5ex]
\textbf{Directores}\\
{Juan Carlos Torres Cantero y Luis López Escudero\\}\\[2cm]
\includegraphics[width=0.3\textwidth]{imagenes/etsiit_logo.png}\\[0.1cm]
\textsc{Escuela Técnica Superior de Ingenierías Informática y de Telecomunicación}\\
\textsc{---}\\
Granada, Septiembre de 2025
\end{minipage}
%\addtolength{\textwidth}{\centeroffset}
%\vspace{\stretch{2}}
\end{titlepage}



\thispagestyle{empty}

\begin{center}
{\large\bfseries Desarrollo software de una aplicación de modelado 3d basada en Signed Distance Funtion (sdf) \\ Motor de renderizado en tiempo real usando Dawn, C++ y WebGPU }\\
\end{center}
\begin{center}
Pablo Cantudo Gómez\\
\end{center}

\vspace{0.7cm}

\vspace{0.5cm}
\noindent\textbf{Palabras clave}: \textit{SDF}, \textit{WebGPU}, \textit{C++}, \textit{ray marching}, \textit{motor de renderizado}
\vspace{0.7cm}

\noindent\textbf{Resumen}\\
\break
Este trabajo presenta el desarrollo de \textit{Copper}, un motor de renderizado 3D en tiempo real implementado en C++ utilizando la API WebGPU a trav\'es de Dawn.
El motor emplea t\'ecnicas de \textit{ray marching} sobre funciones de distancia con signo (SDF), lo que permite una representaci\'on eficiente y flexible de geometr\'ia 
 tridimensional. Una de las principales ventajas de utilizar SDF es la capacidad de expresar de forma directa y compacta operaciones booleanas entre primitivas geom\'etricas
  ---como uniones, intersecciones o diferencias--- mediante simples expresiones algebraicas. Adem\'as, a diferencia de los m\'etodos tradicionales basados en mallas o 
  pol\'igonos, las SDF permiten aplicar operaciones booleanas suaves, como la uni\'on con suavizado (\textit{smooth union}), 
  de forma pr\'acticamente trivial desde el punto de vista computacional.
\bigbreak
Este enfoque simplifica notablemente la construcci\'on de formas complejas, evitando problemas t\'ipicos de la computaci\'on geom\'etrica como el manejo 
de v\'ertices, normales o topolog\'ias complejas. Tambi\'en se facilita la animaci\'on y transformaci\'on de objetos mediante funciones continuas. El motor 
incluye un sistema de sombreado en WGSL, con soporte para sombras suaves, operaciones modulares sobre la escena, seleccion de objetos, carga y guardado de modelos. Los resultados demuestran que
el uso de SDF no solo ofrece un modelo m\'as elegante para definir geometr\'ia, sino que permite construir escenas visualmente complejas con menos c\'odigo y una mayor expresividad gr\'afica.
En conjunto, el sistema demuestra la viabilidad t\'ecnica y creativa del uso de SDF en entornos modernos.
\cleardoublepage

\begin{center}
	{\large\bfseries Development of a 3D modeling application based on Signed Distance Functions (SDF) \\ Real-time rendering engine using Dawn, C++ and WebGPU}\\
\end{center}
\begin{center}
	Pablo Cantudo Gómez\\
\end{center}
\vspace{0.5cm}
\noindent\textbf{Keywords}: \textit{SDF}, \textit{WebGPU}, \textit{C++}, \textit{ray marching}, \textit{rendering engine}
\vspace{0.7cm}

\noindent\textbf{Abstract}\\
\break
This work presents the development of Copper, a real-time 3D rendering engine implemented in C++ using the WebGPU API through Dawn.
The engine employs ray marching techniques over Signed Distance Functions/Fields (SDF), enabling an efficient and flexible representation of three-dimensional geometry.
One of the main advantages of using SDF is the ability to directly and compactly express Boolean operations between geometric primitives ---such as unions, intersections, or differences--- through simple algebraic expressions.
Furthermore, unlike traditional mesh- or polygon-based methods, SDF allows for smooth Boolean operations, such as smooth union, in a virtually trivial manner from a computational standpoint.
\bigbreak
This approach greatly simplifies the construction of complex shapes, avoiding typical problems in geometric computing such as handling vertices, normals, or complex topologies.
It also facilitates the animation and transformation of objects through continuous functions.
The engine includes a shading system in WGSL, with support for soft shadows, modular scene operations, object selection, and model loading and saving.
The results show that the use of SDF not only offers a more elegant model for defining geometry, but also enables the creation of visually complex scenes with less code and greater graphical expressiveness.
Overall, the system demonstrates the technical and creative feasibility of using SDF in modern environments.

\cleardoublepage

\thispagestyle{empty}

\noindent\rule[-1ex]{\textwidth}{2pt}\\[4.5ex]

D. \textbf{Tutora/e(s)}, Profesor(a) del ...

\vspace{0.5cm}

\textbf{Informo:}

\vspace{0.5cm}

Que el presente trabajo, titulado \textit{\textbf{Chief}},
ha sido realizado bajo mi supervisión por \textbf{Estudiante}, y autorizo la defensa de dicho trabajo ante el tribunal
que corresponda.

\vspace{0.5cm}

Y para que conste, expiden y firman el presente informe en Granada a Junio de 2018.

\vspace{1cm}

\textbf{El/la director(a)/es: }

\vspace{5cm}

\noindent \textbf{(nombre completo tutor/a/es)}

\chapter*{Agradecimientos}
Quiero expresar mi agradecimiento a todas las personas que han contribuido directa o indirectamente a la realización de este trabajo.\break

En primer lugar, a mi tutor Juan Carlos Torres Cantero y cotutor Luis López Escudero por la ayuda durante el desarrollo y la correci\'on del proyecto,
a mis amigos por las discusiones y las ideas, y a mi familia por su apoyo incondicional.\break

 Finalmente, me gustaría mencionar a la comunidad de desarrolladores y recursos abiertos, especialmente los foros, artículos y proyectos relacionados con WebGPU, SDF y ray marching,
 cuya documentación y ejemplos han sido una fuente de aprendizaje invaluable.




%\frontmatter
\tableofcontents
\listoffigures
%\listoftables
%
%\mainmatter
%\setlength{\parskip}{5pt}

\chapter{Introducción}
El desarrollo de los gráficos tridimensionales por computadora ha sido un campo de investigación activo y en constante evolución desde sus inicios.\\
La capacidad de crear representaciones visuales de objetos y escenas en tres dimensiones ha revolucionado diversas industrias, desde el entretenimiento hasta la medicina y la ingeniería.\\
A medida que la tecnología avanza, también lo hacen las técnicas y herramientas utilizadas para generar gráficos 3D, lo que plantea nuevos desafíos y oportunidades para los investigadores y desarrolladores.

En sus inicios, la generación de gráficos 3D estaba limitada por la capacidad de cómputo y se basaba en \textit{pipelines} gráficos fijos compuestos por etapas de transformación, iluminación y rasterización.\\
Posteriormente, con la llegada de los \textit{shaders} programables en GPU (Nvidia GeForce 3 en 2001), fue posible sustituir los \textit{pipelines} fijos por \textit{pipelines} programables, lo que abrió un abanico de posibilidades para la creación de efectos visuales complejos y personalizados, como los algoritmos no basados en polígonos.\\
Esto impulsó la investigación en técnicas de representación más avanzadas, como el \textit{ray tracing} y, en particular, el \textit{ray marching}.

En el contexto del renderizado basado en funciones implícitas, las \textit{Signed Distance Functions} (SDF) no son una invención reciente, sino que tienen sus raíces en trabajos mucho más antiguos.\\
El concepto de combinar funciones implícitas mediante operaciones booleanas se remonta al trabajo de Ricci en 1972~\cite{Ricci:1973:CGC}, y fue ampliado en 1989 por B.~Wyvill y G.~Wyvill con el modelado de \textit{soft objects}~\cite{wyvill1989}.\\
Ese mismo año, Sandin, Hart y Kauffman aplicaron \textit{ray marching} a SDF para renderizar fractales tridimensionales~\cite{hart1989ray}.\\
Posteriormente, en 1995, Hart documentó de nuevo la técnica, a la que denominó \textit{Sphere Tracing}~\cite{hart1996}.

La popularización moderna de las SDF en el ámbito del \textit{renderizado en tiempo real} se debe en gran parte a la comunidad \textit{demoscene}, especialmente a partir de mediados de la década de 2000.\\
Trabajos como el de Crane (2005) y Evans (2006) introdujeron la idea de restringir el campo a una distancia euclidiana real, mejorando el rendimiento y la calidad visual.\\
Sin embargo, el uso de esta técnica no está tan extendido en aplicaciones de renderizado en tiempo real, a pesar de su potencial para crear gráficos visuales y eficientes.
\section{Motivación}

Como cualquier persona nacida en los dos mil, he crecido rodeado de videojuegos y el avance en la tecnología de gráficos 3D ha sido un aspecto fascinante de esta industria. 
Durante la carrera de informática, estudié diversas asignaturas relacionadas con gráficos por computadora, cuyos proyectos despertaron y consolidaron mi interés en este campo.

El desarrollo de motores gráficos y técnicas de renderizado siempre me ha resultado un área especialmente atractiva, no solo por su complejidad técnica, 
sino también por el impacto directo que tienen en sectores como el entretenimiento, la simulación o la realidad virtual. A lo largo de mis estudios me encontré con herramientas muy potentes, 
pero también con la dificultad que implica dominarlas o adaptarlas a entornos experimentales. Esto me llevó a plantearme la posibilidad de crear una aplicación propia que sirviera 
como espacio de exploración.

La motivación principal de este trabajo es profundizar en tecnologías emergentes, en concreto \textbf{WebGPU}, un estándar reciente que promete unificar el desarrollo gráfico multiplataforma 
con un acceso eficiente a las GPU modernas. Asimismo, me interesaba experimentar con el modelado mediante \textbf{funciones de distancia (SDF)}, que representan una alternativa flexible al 
modelado poligonal clásico. Considero que la combinación de ambas tecnologías constituye un terreno de investigación con un gran potencial, tanto en aplicaciones prácticas como en entornos 
educativos.

Finalmente, este proyecto me ofrece la oportunidad de afianzar mis conocimientos en programación gráfica, shaders y arquitecturas modernas de GPU, a la vez que desarrollo un software propio 
que pueda servir de base para futuros trabajos de investigación o aplicaciones más complejas en el ámbito del diseño 3D.

\section{Descripción del problema}

El campo del modelado y renderizado 3D ha estado tradicionalmente dominado por herramientas complejas y de gran envergadura, como Blender, Maya o 3ds Max. Si bien estas aplicaciones 
ofrecen una gran potencia y versatilidad, presentan también limitaciones importantes: requieren elevados recursos de hardware, poseen curvas de aprendizaje pronunciadas y no siempre resultan 
adecuadas para entornos de experimentación ligera o proyectos educativos.

Por otro lado, las API gráficas más extendidas, como OpenGL o DirectX, han demostrado su eficacia a lo largo de los años, pero presentan restricciones en cuanto a eficiencia y portabilidad 
en plataformas modernas. El reciente estándar \textbf{WebGPU} surge como respuesta a estas carencias, ofreciendo un modelo de programación más cercano al hardware y multiplataforma, con el 
objetivo de unificar el desarrollo gráfico en navegadores y aplicaciones nativas.

En el ámbito del modelado, el paradigma poligonal sigue siendo el más utilizado, pero alternativas como las \textbf{funciones de distancia (SDF)} permiten representar geometrías complejas 
de manera más compacta y flexible, facilitando la combinación de primitivas y operaciones booleanas. Sin embargo, la integración de estas técnicas en aplicaciones prácticas 
todavía es limitada, especialmente en combinación con tecnologías emergentes como WebGPU.

El problema que aborda este trabajo consiste en la falta de herramientas ligeras que sirvan como demostración y entorno de experimentación para el modelado y renderizado basados en SDF 
sobre WebGPU.

\section{Objetivos}

\subsection*{Objetivo general}
El objetivo principal de este Trabajo de Fin de Grado es el desarrollo de una aplicación de diseño 3D basada en \textbf{WebGPU} y en técnicas de 
\textbf{modelado mediante funciones de distancia (SDF)}, que permita explorar y demostrar el potencial de estas tecnologías como alternativa al modelado poligonal clásico y 
como herramienta de experimentación en el ámbito de los gráficos por computadora.

\subsection*{Objetivos específicos}
Para alcanzar este objetivo general, se plantean los siguientes objetivos específicos:

\begin{itemize}
    \item Investigar y comprender en profundidad el funcionamiento de la API WebGPU y su integración en aplicaciones nativas mediante la librería Dawn.
    \item Diseñar e implementar un motor de renderizado basado en \textit{ray marching} sobre funciones de distancia, capaz de representar primitivas y 
    combinaciones mediante operaciones booleanas y suaves.
    \item Incorporar técnicas de sombreado y efectos visuales (iluminación, sombras suaves) que mejoren la calidad del renderizado.
    \item Desarrollar una interfaz gráfica sencilla que permita al usuario interactuar con la escena y manipular las primitivas.
\end{itemize}


\section{Estructura de la memoria}

La presente memoria se organiza en los siguientes capítulos:

\begin{itemize}
    \item \textbf{Introducción:} Se expone el contexto del trabajo, la motivación, los objetivos planteados y la justificación de la elección de las tecnologías empleadas.
    \item \textbf{Fundamentos teóricos:} Se revisan los conceptos clave de modelado y renderizado 3D, las funciones de distancia (\textit{Signed Distance Functions, SDF}) y el estándar WebGPU, contextualizando el trabajo en el estado actual de la tecnología.
    \item \textbf{Arquitectura de la aplicación:} Se describe la estructura general de la aplicación Copper, detallando los principales módulos, componentes y la interacción entre ellos.
    \item \textbf{Implementación:} Este capítulo se divide en varias secciones:
    \begin{itemize}
        \item \textbf{Ventana y motor de renderizado:} Se explica el proceso de inicialización de la ventana, integración con Dawn/WebGPU y el desarrollo del motor de renderizado basado en SDF y ray marching.
        \item \textbf{Controles e interacción:} Se detalla la implementación de los controles de cámara y de manipulación de objetos mediante la interfaz gráfica.
        \item \textbf{Shaders y efectos visuales:} Se describe el desarrollo de los shaders WGSL utilizados, incluyendo técnicas de sombreado, iluminación y efectos visuales implementados.
    \end{itemize}
    \item \textbf{Conclusiones y trabajos futuros:} Se realiza un balance del trabajo realizado, se revisa el grado de cumplimiento de los objetivos y se plantean posibles líneas de investigación y desarrollo futuras.
    \item \textbf{Bibliografía y anexos:} Se recopilan las referencias bibliográficas consultadas y se incluyen materiales complementarios relevantes para la comprensión y reproducibilidad del trabajo.
\end{itemize}
% \input{Secciones/02_Estado_del_arte}
% \input{Secciones/03_Metodologia}
% \input{Secciones/04_Planificacion}
% \input{Secciones/05_Implementacion}
% \input{Secciones/06_Fuente_recetas}
% \input{Secciones/07_prototipo}
% \chapter{Pruebas}

% \input{Secciones/09_despliegue}
% \input{Secciones/10_costes}
% \input{Secciones/11_conclusiones}

%
\nocite{*}
\addcontentsline{toc}{chapter}{Bibliografía}
\bibliographystyle{miunsrturl}
%\printbibliography
%
\appendix
% \input{apendices/12_Anexo_BD}
% \input{apendices/13_Anexo_prototipo}
% \input{apendices/14_Anexo_pruebas}

\newglossaryentry{nombre en pdf}{
name={Nombre definicion},
description={definicion}
}

\glsaddall
\addcontentsline{toc}{chapter}{Glosario}
\printglossaries
\chapter*{}
\thispagestyle{empty}

\end{document}